\documentclass[a4paper,10pt,twoside]{article}
\usepackage{its_conf}
\begin{document}
\setcounter{section}{0}
\setcounter{figure}{0}
\setcounter{table}{0}
\setcounter{equation}{0}
\setcounter{secnumdepth}{1}
\setcounter{secnumdepth}{1}
% ======= Начало Ваших тезисов. Все, что выше - не изменять ! ===========

\topic{Представления графа диаграммы: память и скорость}

\information{Хаджинова~К.~А.}
{Белорусский государственный университет информатики и радиоэлектроники}
{Минск, Республика Беларусь}
{kseniyakhadzhynava@gmail.com}

% ========= Аннотация ==========
\annotation{Сравниваются объектная модель, списки смежности и типизированные массивы для хранения графа диаграммы на графах до 50\,000 узлов. Массивы занимают в десять раз меньше памяти и обходят граф в десять раз быстрее, колоночный JSON на 41\,\% компактнее строчного.}

\begin{multicols}{2} %НЕ УДАЛЯТЬ!!!

\section*{Введение}
В редакторах диаграмм модель процесса, данных или архитектуры хранится как граф: узлы соответствуют элементам, ребра~-- связям. При малом числе узлов способ хранения не важен, но при числе узлов порядка $10^4$ и более он определяет объем памяти, время открытия документа и скорость обхода графа при проверке и отрисовке.

Для больших графов важны компактность хранения и скорость обработки данных [1, 2]; обзор методов работы с графами в информационной визуализации приведен в [3]. Для хранения графов применяют объекты со ссылками, списки смежности и компактные массивы [4]. Объектная модель удобна для редактирования, но каждый узел и ребро несут накладные расходы объекта.

Цель работы~-- измерить, как выбор представления влияет на память, время построения и обхода, а также на размер сериализации. Ранее авторами была предложена конвейерная модель обработки диаграмм с единым графовым представлением [5]; здесь исследуется его хранение.

\section{Представления графа}
Рассматриваются три представления одного и того же графа. В объектной модели узел и ребро~-- объекты, которые ссылаются друг на друга. В списках смежности узлы хранятся в массиве, а для каждого узла ведутся массивы номеров соседей. В представлении на типизированных массивах координаты узлов и концы ребер лежат в непрерывных массивах, а смежность задана смещениями.

Выигрыш по памяти оценивается отношением по формуле~(1):
$$
K_M = M_{obj}/M_{arr},\eqno(1)
$$
где $M_{obj}$ и $M_{arr}$~-- память объектной модели и массивов. Для JSON сравниваются строчный формат (массив объектов) и колоночный (по массиву на каждое поле).

\section{Методика эксперимента}
Измерения выполнены на отдельном стенде (Node.js 24, движок V8; Intel Core i5-1035G1, 24~ГБ ОЗУ, Windows~11). Генератор строит текст на языке диаграмм для графов из 500, 2\,000, 10\,000 и 50\,000 узлов (ребер в полтора раза больше) трех типов: случайных, иерархических и с узлами-хабами; начальное значение фиксировано. Измеряются память модели (куча после сборки мусора), время построения и полного обхода, размер JSON, время его записи и чтения. Каждое измерение выполняется в отдельном процессе, 11 запусков, первый отбрасывается; приводятся медианы, для трех типов графов~-- медиана по типам.

\section{Результаты}
Разбор текста занимает около 0,55~с на 50\,000 узлов и не зависит от представления, поэтому в табл.~\ref{tab:repr} он не учтен. Память растет пропорционально числу узлов (рис.~\ref{fig:mem}); по формуле (1) $K_M$ равен 10,3 при 10\,000 узлов и 10,6 при 50\,000. В пересчете на один узел с его ребрами это около 510 байт для объектной модели и 48 байт для массивов.

\myTable{Представления графа\label{tab:repr}}{{\columnwidth}{|>{\raggedright\arraybackslash}p{2.5cm}|X|X|X|} \hline
Показатель & Объекты & Списки & Массивы \\ \hline
Память, МБ (10\,000) & 4,87 & 4,64 & 0,47 \\ \hline
Память, МБ (50\,000) & 24,3 & 23,2 & 2,30 \\ \hline
Построение, мс (10\,000) & 21,4 & 13,4 & 6,6 \\ \hline
Построение, мс (50\,000) & 200 & 98 & 31 \\ \hline
Обход, мс (10\,000) & 18,3 & 7,5 & 3,2 \\ \hline
Обход, мс (50\,000) & 138 & 27 & 14 \\ \hline
}

Списки смежности экономят лишь 5\,\% памяти, но обходят граф в 5 раз быстрее объектов, а массивы~-- в 10 раз.

\image{
[width=0.8\columnwidth]{fig-memory-ru.png}
\caption{Память внутренней модели графа\label{fig:mem}}
}

Колоночный JSON занимает на 41\,\% меньше места, записывается в 2,1 раза и читается в 5,2 раза быстрее строчного (табл.~\ref{tab:json}); чтение включает построение соответствующего представления.

\myTable{Сериализация, 50\,000 узлов\label{tab:json}}{{\columnwidth}{|>{\raggedright\arraybackslash}p{2.1cm}|X|X|X|} \hline
Формат & Размер, КБ & Запись, мс & Чтение, мс \\ \hline
Строчный & 4482 & 29,0 & 43,3 \\ \hline
Колоночный & 2626 & 13,6 & 8,3 \\ \hline
}

Размер файла в обоих форматах растет линейно с числом узлов (рис.~\ref{fig:json}), при этом колоночный формат остается компактнее на всех размерах.

\image{
[width=0.8\columnwidth]{fig-json-ru.png}
\caption{Размер JSON-файла графа\label{fig:json}}
}

При росте числа узлов в 100 раз время обхода возрастает в 136 раз для объектной модели, в 39 раз для списков смежности и в 26 раз для массивов (рис.~\ref{fig:trav}). Причина различий, вероятно, состоит в расположении данных в памяти [6] и требует отдельной проверки.

\image{
[width=0.8\columnwidth]{fig-trav-ru.png}
\caption{Время полного обхода графа\label{fig:trav}}
}

Разбор текста составляет 75\,\% времени загрузки объектной модели, 85\,\% для списков и 95\,\% для массивов при 50\,000 узлов (рис.~\ref{fig:load}). Поэтому выбор представления ускоряет загрузку лишь частично, а основной резерв дает ускорение разбора.

\image{
[width=0.8\columnwidth]{fig-load-ru.png}
\caption{Структура времени загрузки, 50\,000 узлов\label{fig:load}}
}

Стоимость добавления и удаления ребер не измерялась; для массивов она выше, поэтому они подходят для хранения и анализа, а объекты~-- для интерактивного редактирования.

\section*{Заключение}
Типизированные массивы сокращают память в десять раз и ускоряют построение и обход, а колоночный JSON уменьшает размер файла и ускоряет загрузку. При 500 узлах память модели не превышает 0,3~МБ для любого представления, поэтому выбор представления на малых графах не влияет на результат. Ограничение работы~-- измерения выполнены в Node.js, а операции редактирования не рассматривались.

\ListReferences{
\item Козарь,~Р.~В. Методы кластеризации данных при распознавании эндоскопических снимков в задачах компьютерной медицинской диагностики~/ Р.~В.~Козарь, Н.~С.~Конойко, А.~А.~Навроцкий~// Доклады БГУИР.~-- 2023.~-- Т.~21, №~1.~-- С.~94--97.~-- DOI: 10.35596/1729-7648-2023-21-1-94-97.
\item Abello,~J. ASK-GraphView: a large scale graph visualization system~/ J.~Abello, F.~van Ham, N.~Krishnan~// IEEE Trans. Vis. Comput. Graph.~-- 2006.~-- Vol.~12, №~5.~-- P.~669--676.
\item Herman,~I. Graph visualization and navigation in information visualization: a survey~/ I.~Herman, G.~Melancon, M.~S.~Marshall~// IEEE Trans. Vis. Comput. Graph.~-- 2000.~-- Vol.~6, №~1.~-- P.~24--43.
\item Introduction to Algorithms~/ T.~H.~Cormen [et al.].~-- 4th ed.~-- Cambridge: MIT Press, 2022.~-- 1312~p.
\item Ломако,~А.~В. Автоматизация преобразования изображений диаграмм в программный код~/ А.~В.~Ломако, К.~А.~Хаджинова~// Технологии передачи и обработки информации.~-- 2026.
\item Drepper,~U. What every programmer should know about memory.~-- Red Hat, 2007.~-- 114~p.
}

\end{multicols} % НЕ УДАЛЯТЬ!

% ===== Конец материалов. Все, что ниже, не изменять !===========
\end{document}
